\begin{itemframe}{صورت بندی برنامه‌ریزی خطی}
\itm
برنامه‌ریزی خطی کاربردهای بسیاری دارد. در این بخش، بررسی می‌کنیم که چگونه چند مسائل را می‌توان به‌صورت برنامه‌ریزی خطی فرمول‌بندی کرد.
\itm
در بحث برنامه ریزی خطی، مهم‌ترین مهارت تشخیص این است که چه زمانی می‌توان مسئله‌ای را به‌صورت یک برنامه‌ی خطی نوشت.
اگر مسئله‌ای را به برنامه‌ای خطی با اندازه‌ی چندجمله‌ای تبدیل کنید، می‌توان آن را در زمان چندجمله‌ای حل کرد.
بسته‌های نرم‌افزاری متعددی برای حل برنامه‌های خطی به صورت کارآمد وجود دارند.
\end{itemframe}


\begin{itemframe}{صورت بندی برنامه‌ریزی خطی}
\itm
پیش‌تر هنگام حل مسائل گراف از نشانه گذاری‌هایی مانند
$v:d$
و
$(u,v):f$
استفاده می‌کردیم.

با این حال در برنامه‌های خطی، هنگام بیان متغیرها، رأس‌ها و یال‌ها را با زیرنویس مشخص می‌کنیم.
مثلاً وزن کوتاه‌ترین مسیر برای رأس v را به صورت
$d_v$
می‌نویسیم.
\itm
مسائلی مانند کوتاه‌ترین مسیر تک مبداء و شار بیشینه، با استفاده از برنامه‌ریزی خطی برای قابل حل اند. اما برای این مسائل قبلاً نیز الگوریتم‌های کارآمدی می‌شناختیم.
\itm
در واقع، یک الگوریتم که برای مسئله‌ی خاصی طراحی شده باشد، اغلب کاراتر از برنامه‌ریزی خطی است. (مانند الگوریتم دیکسترا برای مسئله‌ی کوتاه‌ترین مسیر تک‌مبدأ)

\end{itemframe}


\begin{itemframe}{صورت بندی برنامه‌ریزی خطی}
\itm
اما قدرت واقعی برنامه‌ریزی خطی از توانایی آن در حل مسائل جدید ناشی می‌شود.
\itm
برای مثال مسئله انتخابات را که در ابتدای این فصل با آن روبه‌رو شدیم در نظر بگیرید. این مسئله توسط هیچ‌یک از الگوریتم‌هایی که تا کنون مطالعه کرده‌ایم حل نمی‌شود، اما با برنامه‌ریزی خطی می‌توان آن را حل کرد.
\itm
برنامه‌ریزی خطی به‌ویژه زمانی مفید است که بخواهیم گونه‌های تعمیم‌یافته‌ی مسائلی را حل کنیم که هنوز الگوریتم کارآمدی برای آن‌ها نمی‌شناسیم. به‌عنوان نمونه، تعمیم‌هایی از مسئله‌ی شار بیشینه را بررسی خواهیم کرد.
\end{itemframe}


\begin{itemframe-s}{صورت بندی برنامه‌ریزی خطی}{شار با کمترین هزینه}
\itm
فرض کنید در مسئله شار بیشینه، علاوه بر یک ظرفیت
$c(u,v)$
برای هر یال
$(u,v)$,
یک هزینه‌ی حقیقی

$a(u,v)$
نیز داده شده است.
\itm
به طوری که اگر
$f_{uv}$
واحد شار از یال
$(u,v)$
ارسال کنید، هزینه‌ای برابر
$a(u,v) \cdot f_{uv}$
خواهد داشت.
\itm
هدف آن است که دقیقاً d واحد شار از s به t ارسال کنیم، به‌گونه‌ای که مجموع هزینه
$$
\sum_{(u,v) \in E} a(u,v) \cdot f_{uv}
$$
کمینه شود. این مسئله به نام مسئله‌ی شار با کمترین هزینه
\fn{minimum-cost fow problem}
 شناخته می‌شود.
\end{itemframe-s}


\begin{itemframe-s}{صورت بندی برنامه‌ریزی خطی}{شار با کمترین هزینه}
\itm
شکل زیر نمونه‌ای از این مسئله را نشان می‌دهد. هدف ارسال ۴ واحد شار از s به t با کمترین هزینه است.
\pic{3}


\end{itemframe-s}


\begin{itemframe-s}{صورت بندی برنامه‌ریزی خطی}{شار با کمترین هزینه}
\itm
هر شار مجاز، هزینه‌ای برابر
$$
\sum_{(u,v) \in E} a(u,v) \cdot f_{uv}
$$
تحمیل می‌کند.
\itm
پرسش این است: کدام شار ۴ واحدی این هزینه را کمینه می‌کند؟
\end{itemframe-s}


\begin{itemframe-s}{صورت بندی برنامه‌ریزی خطی}{شار با کمترین هزینه}
\itm
 شکل زیر یک راه‌حل بهینه را نشان می‌دهد که هزینه‌ی کل آن نوشته ‌شده است:

\pic{4}
$$
\sum_{(u,v)\in E} a(u,v) \cdot f_{uv}
= (2 \cdot 2) + (5 \cdot 2) + (3 \cdot 1) + (7 \cdot 1) + (1 \cdot 3) = 27
$$


\end{itemframe-s}


\begin{itemframe-s}{صورت بندی برنامه‌ریزی خطی}{شار با کمترین هزینه}
\itm
الگوریتم‌هایی با زمان چندجمله‌ای که به‌طور خاص برای حل این مسئله طراحی شده‌اند وجود دارند. با این حال، مسئله‌ی شار با کمترین هزینه را می‌توان به‌صورت یک برنامه‌ی خطی بیان کرد:

\begin{align*}
\text{minimize} \quad
& \sum_{(u,v) \in E} a(u,v) \cdot f_{uv} \\
\text{to subject} \quad \\
& f_{uv} \leq c(u,v) \quad \text{ each for} u,v \in V \\
& \sum_{v \in V} f_{vu} - \sum_{v \in V} f_{uv} = 0
    \quad \text{ each for} u \in V - \{s,t\} \\
& \sum_{v \in V} f_{sv} - \sum_{v \in V} f_{vs} = d \\
& f_{uv} \geq 0 \quad \text{ each for} u,v \in V
\end{align*}


\end{itemframe-s}


\begin{itemframe-s}{صورت بندی برنامه‌ریزی خطی}{شار چندکالایی}
\itm
فرض کنید یک کارخانه توپ‌های هاکی تولید می‌کند و باید آن‌ها را به از انبار برساند.
این وضعیت به‌صورت یک شبکه شار مدل می‌شود: گره مبدأ نمایانگر محل تولید، گره‌های میانی مسیرهای حمل‌ونقل، و گره مقصد نمایانگر انبار است. در مسئله شار بیشینه هدف این بود که مشخص کنیم حداکثر چه تعداد توپ می‌تواند از کارخانه به انبار برسد.
\itm
فرض کنید در بخش جدید، اقلام دیگری مثل چوب و کلاه هاکی هم تولید می‌شود. هر قطعه در کارخانه‌ی مخصوص به خود ساخته می‌شود، انبار جداگانه‌ای دارد، و هر روز باید از کارخانه به انبار منتقل شود. این نمونه‌ای از مسئله‌ی شار چندکالایی
\fn{multi-commodity flow}
 است.
\end{itemframe-s}


\begin{itemframe-s}{صورت بندی برنامه‌ریزی خطی}{شار چندکالایی}
\itm
فرض کنید، k کالای مختلف داریم:
$$
K_1, K_2, \dots, K_k
$$
که هر کالای i با سه‌تایی
$$
K_i = (s_i, t_i, d_i)
$$
تعریف می‌شود. که در آن، رأس $s_i$ مبدأ کالای i، رأس $ t_i$ مقصد آن، و $d_i$ شار مورد نیاز آن کالا است.

\end{itemframe-s}


\begin{itemframe-s}{صورت بندی برنامه‌ریزی خطی}{شار چندکالایی}
\itm
برای هر کالا i شاری به‌صورت تابع حقیقی $f_i$ تعریف می‌شود، به‌گونه‌ای که $f_{iuv}$ شار کالای i از رأس u به رأس v باشد.
این تابع باید قیود ظرفیت و پایستگی شار را رعایت کند. شار کل
\fn{aggregate flow}
 روی یال (u,v) برابر است با مجموع شار‌های کالاهای مختلف:
$$
f_{uv} = \sum_{i=1}^k f_{iuv}
$$

و باید قید
$$
f_{uv} \leq c(u,v)
$$

را رعایت کند.
\itm
این مسئله تابع هدفی ندارد؛ پرسش فقط این است که آیا چنین شاری وجود دارد یا نه. بنابراین برنامه‌ی خطی این مسئله دارای تابع هدف تهی
\fn{null objective function}
 است:
\end{itemframe-s}


\begin{itemframe-s}{صورت بندی برنامه‌ریزی خطی}{شار چندکالایی}
\itm
 صورت بندی مسئله شار چند کالایی به این شکل است:
\begin{align*}
\text{minimize} \quad
& 0 \\[3pt]
\text{ to subject} \quad
& \sum_{i=1}^k f_{iuv} \leq c(u,v)
   \quad \text{ each for }u,v \in V \\[3pt]
& \sum_{v \in V} f_{iuv} - \sum_{v \in V} f_{ivu} = 0
   \quad \text{ each for } i = 1,2,\dots,k  \text{ each and }
   u \in V - \{s_i,t_i\} \\[3pt]
& \sum_{v \in V} f_{is_iv} - \sum_{v \in V} f_{iv s_i} = d_i
   \quad \text{ each for } i = 1,2,\dots,k \\[3pt]
& f_{iuv} \geq 0
   \quad \text{ each for } u,v \in V \text{ each and }
   i = 1,2,\dots,k
\end{align*}

\end{itemframe-s}


\begin{itemframe-s}{صورت بندی برنامه‌ریزی خطی}{شار چندکالایی}
\itm
جالب است بدانید تنها الگوریتم چندجمله‌ای شناخته‌شده برای این مسئله، بر مبنای برنامه‌ ریزی خطی است.

\end{itemframe-s}